% ============================================================
% 05_METODOLOGIA.TEX
% Sistema Adaptativo de Classificação Multicritério
% de Fluidos de Corte
%
% FINALIDADE
%
% Apresentar a metodologia multicritério prevista para o
% projeto:
%
% 1. definição dos critérios;
% 2. construção da matriz de decisão;
% 3. normalização;
% 4. ponderação pelo método CRITIC;
% 5. classificação pelo TOPSIS;
% 6. cenários;
% 7. análise de Pareto;
% 8. sensibilidade;
% 9. robustez.
%
% Este arquivo constitui apenas a BASE da apresentação.
%
% Os integrantes deverão posteriormente preencher:
% - critérios realmente utilizados;
% - normalização escolhida;
% - parâmetros;
% - cenários;
% - resultados.
%
% ============================================================


% ============================================================
% FRAME 1 — VISÃO GERAL
% ============================================================

\begin{frame}{Metodologia}

    \begin{block}{Pipeline multicritério}
        \centering

        Base canônica
        $\rightarrow$
        Critérios
        $\rightarrow$
        Normalização
        $\rightarrow$
        CRITIC
        $\rightarrow$
        TOPSIS
        $\rightarrow$
        Robustez
    \end{block}

    \vspace{0.4cm}

    O sistema foi estruturado para separar três questões:

    \begin{itemize}
        \item \textbf{o que será comparado}: critérios;
        \item \textbf{quanto cada critério informa}: CRITIC;
        \item \textbf{como as alternativas serão ordenadas}: TOPSIS.
    \end{itemize}

    \vspace{0.3cm}

    Análises de Pareto, cenários e sensibilidade serão utilizadas
    posteriormente para avaliar a estabilidade da classificação.

\end{frame}


% ============================================================
% FRAME 2 — DEFINIÇÃO DOS CRITÉRIOS
% ============================================================

\begin{frame}{Definição dos critérios}

    Os critérios não deverão ser escolhidos apenas porque existem
    na planilha.

    \vspace{0.3cm}

    A seleção deverá considerar:

    \begin{itemize}
        \item significado técnico;
        \item qualidade dos dados;
        \item capacidade discriminatória;
        \item redundância;
        \item disponibilidade;
        \item interpretação;
        \item direção de preferência.
    \end{itemize}

    \vspace{0.3cm}

    \begin{block}{Resultado esperado}
        Uma lista de critérios justificável estatística e
        tecnicamente.
    \end{block}

\end{frame}


% ============================================================
% FRAME 3 — TIPOS DE CRITÉRIOS
% ============================================================

\begin{frame}{Direção dos critérios}

    Cada variável deverá possuir uma regra de preferência conhecida.

    \vspace{0.3cm}

    \begin{columns}[T]

        \begin{column}{0.31\textwidth}

            \begin{block}{BENEFIT}
                Quanto maior o valor, melhor a alternativa.
            \end{block}

        \end{column}


        \begin{column}{0.31\textwidth}

            \begin{block}{COST}
                Quanto menor o valor, melhor a alternativa.
            \end{block}

        \end{column}


        \begin{column}{0.31\textwidth}

            \begin{block}{TARGET}
                O melhor desempenho ocorre próximo de um valor-alvo.
            \end{block}

        \end{column}

    \end{columns}

    \vspace{0.4cm}

    Além disso, algumas variáveis poderão ser utilizadas apenas como:

    \begin{center}
        \textbf{restrição}
        \qquad ou \qquad
        \textbf{diagnóstico}.
    \end{center}

\end{frame}


% ============================================================
% FRAME 4 — CONFORMIDADE
% ============================================================

\begin{frame}{Conformidade e desempenho}

    \begin{block}{Separação importante}
        Conformidade técnica e desempenho multicritério representam
        conceitos diferentes.
    \end{block}

    \vspace{0.4cm}

    Uma alternativa poderá:

    \begin{itemize}
        \item apresentar bom desempenho relativo;
        \item possuir score elevado;
        \item e ainda assim apresentar restrição técnica.
    \end{itemize}

    \vspace{0.3cm}

    Portanto, quando existirem especificações confirmadas, elas
    deverão ser tratadas explicitamente como:

    \begin{itemize}
        \item restrições rígidas;
        \item restrições flexíveis;
        \item ou informações diagnósticas.
    \end{itemize}

\end{frame}


% ============================================================
% FRAME 5 — MATRIZ DE DECISÃO
% ============================================================

\begin{frame}{Matriz de decisão}

    Após a definição dos critérios será construída a matriz:

    \[
    X=[x_{ij}]
    \]

    em que:

    \begin{itemize}
        \item \(i\) representa cada alternativa de fluido;
        \item \(j\) representa cada critério;
        \item \(x_{ij}\) representa o desempenho da alternativa
              \(i\) no critério \(j\).
    \end{itemize}

    \vspace{0.3cm}

    \begin{block}{Antes da utilização}
        A matriz deverá ser verificada quanto a missing, unidades,
        critérios constantes, direções e consistência das alternativas.
    \end{block}

    \vspace{0.2cm}

    \begin{center}
        Número de alternativas:
        \texttt{TO\_BE\_FILLED}

        \vspace{0.1cm}

        Número de critérios:
        \texttt{TO\_BE\_FILLED}
    \end{center}

\end{frame}


% ============================================================
% FRAME 6 — NORMALIZAÇÃO
% ============================================================

\begin{frame}{Normalização}

    Os critérios poderão possuir unidades e escalas diferentes.

    \vspace{0.3cm}

    Por isso será necessária uma transformação para uma escala
    comparável:

    \[
    X=[x_{ij}]
    \quad\longrightarrow\quad
    R=[r_{ij}]
    \]

    \vspace{0.3cm}

    \begin{block}{Objetivo}
        Tornar os critérios comparáveis preservando sua interpretação
        de preferência.
    \end{block}

    \vspace{0.3cm}

    O método principal será definido pelos integrantes durante a
    implementação:

    \begin{center}
        \texttt{PRIMARY\_NORMALIZATION = TO\_BE\_FILLED}
    \end{center}

\end{frame}


% ============================================================
% FRAME 7 — EXEMPLO DE NORMALIZAÇÃO
% ============================================================

\begin{frame}{Exemplo de normalização}

    Uma possibilidade para critérios de benefício é a transformação
    Min--Max:

    \[
    r_{ij}
    =
    \frac{x_{ij}-\min_i(x_{ij})}
         {\max_i(x_{ij})-\min_i(x_{ij})}.
    \]

    Para critérios de custo, uma formulação correspondente é:

    \[
    r_{ij}
    =
    \frac{\max_i(x_{ij})-x_{ij}}
         {\max_i(x_{ij})-\min_i(x_{ij})}.
    \]

    \vspace{0.3cm}

    \begin{alertblock}{Importante}
        Essas equações representam uma possibilidade metodológica.
        A normalização oficial somente será definida durante a execução.
    \end{alertblock}

\end{frame}


% ============================================================
% FRAME 8 — CRITIC
% ============================================================

\begin{frame}{Ponderação objetiva: CRITIC}

    O método CRITIC será utilizado para estimar pesos objetivos dos
    critérios.

    \vspace{0.3cm}

    Ele considera simultaneamente:

    \begin{columns}[T]

        \begin{column}{0.47\textwidth}

            \begin{block}{Contraste}
                A variabilidade apresentada por cada critério.
            \end{block}

        \end{column}


        \begin{column}{0.47\textwidth}

            \begin{block}{Conflito}
                A relação do critério com os demais critérios.
            \end{block}

        \end{column}

    \end{columns}

    \vspace{0.4cm}

    Assim, o peso não é simplesmente definido pela opinião do grupo.

\end{frame}


% ============================================================
% FRAME 9 — FÓRMULA CRITIC
% ============================================================

\begin{frame}{Cálculo do CRITIC}

    Para cada critério \(j\), o conteúdo informacional será calculado por:

    \[
    C_j
    =
    \sigma_j
    \sum_{k=1}^{p}
    (1-r_{jk}),
    \]

    onde:

    \begin{itemize}
        \item \(\sigma_j\) representa a dispersão do critério;
        \item \(r_{jk}\) representa a correlação entre os critérios
              \(j\) e \(k\).
    \end{itemize}

    Os pesos serão então:

    \[
    w_j
    =
    \frac{C_j}
         {\sum_{j=1}^{p}C_j}.
    \]

    \vspace{0.2cm}

    Consequentemente:

    \[
    \sum_{j=1}^{p}w_j=1.
    \]

\end{frame}


% ============================================================
% FRAME 10 — INTERPRETAÇÃO DO CRITIC
% ============================================================

\begin{frame}{Como interpretar os pesos CRITIC?}

    \begin{block}{Maior peso}
        Representa maior conteúdo informacional relativo dentro da
        configuração analisada.
    \end{block}

    \vspace{0.3cm}

    O peso depende da combinação entre:

    \[
    \text{variabilidade}
    +
    \text{conflito com outros critérios}.
    \]

    \vspace{0.3cm}

    \begin{alertblock}{Cuidado}
        Um peso CRITIC elevado não significa automaticamente que o
        critério seja tecnicamente mais importante.
    \end{alertblock}

\end{frame}


% ============================================================
% FRAME 11 — CORRELAÇÃO NO CRITIC
% ============================================================

\begin{frame}{Conflito entre critérios no CRITIC}

    A formulação principal utiliza:

    \[
    1-r_{jk}.
    \]

    \vspace{0.3cm}

    Portanto:

    \begin{itemize}
        \item correlação positiva próxima de \(1\)
              reduz o termo de conflito;

        \item correlações menores aumentam o conflito;

        \item correlação negativa pode aumentar ainda mais esse termo.
    \end{itemize}

    \vspace{0.3cm}

    \begin{alertblock}{Regra metodológica}
        A substituição por \(1-|r_{jk}|\) alteraria o método e não
        deverá ser realizada silenciosamente.
    \end{alertblock}

\end{frame}


% ============================================================
% FRAME 12 — DO CRITIC AO TOPSIS
% ============================================================

\begin{frame}{Do peso à classificação}

    Após o CRITIC teremos:

    \[
    \mathbf{w}
    =
    (w_1,w_2,\ldots,w_p).
    \]

    Esses pesos serão combinados com a matriz normalizada:

    \[
    v_{ij}
    =
    w_j r_{ij}.
    \]

    \vspace{0.3cm}

    A matriz:

    \[
    V=[v_{ij}]
    \]

    será utilizada pelo TOPSIS para comparar cada alternativa com
    referências ideais.

\end{frame}


% ============================================================
% FRAME 13 — TOPSIS
% ============================================================

\begin{frame}{Classificação: TOPSIS}

    O TOPSIS procura alternativas que estejam:

    \begin{center}
        \textbf{próximas da solução ideal}
    \end{center}

    e simultaneamente:

    \begin{center}
        \textbf{distantes da solução anti-ideal}.
    \end{center}

    \vspace{0.3cm}

    O processo envolve:

    \begin{enumerate}
        \item matriz normalizada ponderada;
        \item solução ideal positiva;
        \item solução ideal negativa;
        \item cálculo das distâncias;
        \item coeficiente de proximidade.
    \end{enumerate}

\end{frame}


% ============================================================
% FRAME 14 — IDEAL E ANTI-IDEAL
% ============================================================

\begin{frame}{Soluções ideal e anti-ideal}

    Se todos os critérios estiverem orientados de forma que:

    \[
    \text{maior valor}=\text{melhor desempenho},
    \]

    teremos:

    \[
    A^+
    =
    \left\{
    \max_i v_{ij}
    \right\}
    \]

    e:

    \[
    A^-
    =
    \left\{
    \min_i v_{ij}
    \right\}.
    \]

    \vspace{0.3cm}

    \begin{alertblock}{Cuidado}
        Critérios de custo não devem ser invertidos novamente caso já
        tenham sido orientados durante a normalização.
    \end{alertblock}

\end{frame}


% ============================================================
% FRAME 15 — DISTÂNCIAS TOPSIS
% ============================================================

\begin{frame}{Distâncias às soluções de referência}

    Para cada alternativa \(i\):

    \[
    S_i^+
    =
    \sqrt{
    \sum_j
    (v_{ij}-v_j^+)^2
    }
    \]

    representa sua distância à solução ideal.

    \vspace{0.3cm}

    E:

    \[
    S_i^-
    =
    \sqrt{
    \sum_j
    (v_{ij}-v_j^-)^2
    }
    \]

    representa sua distância à solução anti-ideal.

\end{frame}


% ============================================================
% FRAME 16 — COEFICIENTE TOPSIS
% ============================================================

\begin{frame}{Coeficiente de proximidade TOPSIS}

    O coeficiente de proximidade será:

    \[
    C_i
    =
    \frac{S_i^-}
         {S_i^+ + S_i^-}.
    \]

    Para uma execução válida:

    \[
    0\leq C_i\leq1.
    \]

    \vspace{0.3cm}

    Quanto maior \(C_i\), maior a proximidade relativa da alternativa
    à solução ideal.

    \vspace{0.3cm}

    \begin{alertblock}{Interpretação}
        \(C_i\) não representa a probabilidade de o fluido ser o melhor.
    \end{alertblock}

\end{frame}


% ============================================================
% FRAME 17 — RANKING
% ============================================================

\begin{frame}{Construção do ranking}

    O ranking principal será obtido pela ordenação decrescente de:

    \[
    C_i.
    \]

    \vspace{0.3cm}

    Além da posição, serão avaliados:

    \begin{itemize}
        \item diferença entre alternativas;
        \item proximidade entre primeiro e segundo;
        \item possíveis quase empates;
        \item conformidade;
        \item perfil de forças e fraquezas.
    \end{itemize}

    \vspace{0.3cm}

    \begin{block}{Objetivo}
        Produzir um ranking explicável e não apenas uma sequência de
        posições.
    \end{block}

\end{frame}


% ============================================================
% FRAME 18 — CENÁRIOS
% ============================================================

\begin{frame}{Análise de cenários}

    O cenário principal utilizará a configuração metodológica aprovada.

    \vspace{0.3cm}

    Cenários alternativos poderão avaliar diferentes hipóteses
    decisórias, quando justificadas.

    Exemplos possíveis:

    \begin{itemize}
        \item pesos CRITIC;
        \item pesos iguais;
        \item maior ênfase em determinado grupo de critérios;
        \item diferentes prioridades técnicas.
    \end{itemize}

    \vspace{0.3cm}

    \begin{alertblock}{Regra}
        Os cenários não deverão ser escolhidos com o objetivo de fazer
        uma determinada alternativa vencer.
    \end{alertblock}

\end{frame}


% ============================================================
% FRAME 19 — PARETO
% ============================================================

\begin{frame}{Análise de Pareto}

    A dominância de Pareto fornecerá uma avaliação complementar que
    não depende dos pesos CRITIC.

    \vspace{0.3cm}

    Uma alternativa \(A\) domina \(B\) quando:

    \begin{itemize}
        \item não é pior que \(B\) em nenhum critério;
        \item é estritamente melhor em pelo menos um critério.
    \end{itemize}

    \vspace{0.3cm}

    Alternativas não dominadas formarão a:

    \begin{center}
        \textbf{fronteira de Pareto}.
    \end{center}

    \vspace{0.2cm}

    \begin{alertblock}{Importante}
        Pertencer à fronteira de Pareto não significa automaticamente
        ocupar a primeira posição.
    \end{alertblock}

\end{frame}


% ============================================================
% FRAME 20 — SENSIBILIDADE
% ============================================================

\begin{frame}{Análise de sensibilidade}

    Após o ranking principal será avaliado quanto as conclusões
    dependem das escolhas metodológicas.

    \vspace{0.3cm}

    A análise poderá incluir:

    \begin{itemize}
        \item perturbação dos pesos;
        \item retirada de um critério;
        \item normalizações alternativas;
        \item retirada de uma alternativa;
        \item cenários;
        \item mudança da ordem do ranking.
    \end{itemize}

    \vspace{0.3cm}

    \begin{block}{Pergunta central}
        Pequenas ou razoáveis mudanças na configuração alteram a
        conclusão?
    \end{block}

\end{frame}


% ============================================================
% FRAME 21 — PERTURBAÇÃO DOS PESOS
% ============================================================

\begin{frame}{Sensibilidade aos pesos}

    Poderão ser avaliadas perturbações planejadas, por exemplo:

    \begin{center}
        \(\pm5\%\)
        \qquad
        \(\pm10\%\)
        \qquad
        \(\pm20\%\)
    \end{center}

    \vspace{0.3cm}

    Após cada perturbação, os pesos deverão permanecer:

    \[
    w_j\geq0
    \]

    e:

    \[
    \sum_j w_j=1.
    \]

    \vspace{0.3cm}

    Os integrantes deverão definir a regra definitiva durante a
    implementação.

\end{frame}


% ============================================================
% FRAME 22 — NORMALIZAÇÃO E CRITIC
% ============================================================

\begin{frame}{Sensibilidade à normalização}

    Alterar a normalização pode modificar:

    \begin{itemize}
        \item dispersão dos critérios;
        \item correlações;
        \item pesos CRITIC;
        \item distâncias TOPSIS;
        \item ranking.
    \end{itemize}

    \vspace{0.3cm}

    Portanto, uma comparação metodologicamente consistente deverá
    seguir:

    \begin{block}{Reexecução}
        \centering

        Normalização alternativa
        $\rightarrow$
        novo CRITIC
        $\rightarrow$
        novo TOPSIS
        $\rightarrow$
        comparação
    \end{block}

\end{frame}


% ============================================================
% FRAME 23 — REMOÇÃO DE ALTERNATIVA
% ============================================================

\begin{frame}{Sensibilidade ao conjunto de alternativas}

    Também poderá ser avaliado o efeito de retirar uma alternativa
    por vez.

    \vspace{0.3cm}

    Para cada remoção:

    \begin{block}{Leave-one-alternative-out}
        \centering

        Remover
        $\rightarrow$
        normalizar novamente
        $\rightarrow$
        recalcular CRITIC
        $\rightarrow$
        recalcular TOPSIS
    \end{block}

    \vspace{0.4cm}

    Isso permitirá investigar possíveis alterações na ordem relativa
    das alternativas restantes.

\end{frame}


% ============================================================
% FRAME 24 — RANK REVERSAL
% ============================================================

\begin{frame}{Rank reversal}

    Uma mudança na composição do conjunto pode provocar mudança na
    ordem relativa das alternativas.

    \vspace{0.3cm}

    Esse fenômeno será registrado como:

    \begin{center}
        \textbf{rank reversal}.
    \end{center}

    \vspace{0.3cm}

    \begin{alertblock}{Interpretação}
        A ocorrência de rank reversal não deve ser tratada
        automaticamente como erro. Ela deverá ser analisada dentro da
        natureza relativa do método.
    \end{alertblock}

\end{frame}


% ============================================================
% FRAME 25 — ROBUSTEZ
% ============================================================

\begin{frame}{Avaliação de robustez}

    A robustez sintetizará os resultados das diferentes verificações.

    \vspace{0.3cm}

    Serão observados, conforme a implementação:

    \begin{itemize}
        \item estabilidade do primeiro colocado;
        \item estabilidade do top 3;
        \item mudança de posição;
        \item influência dos critérios;
        \item influência das alternativas;
        \item sensibilidade à normalização;
        \item comportamento nos cenários;
        \item consistência com Pareto.
    \end{itemize}

\end{frame}


% ============================================================
% FRAME 26 — NÃO CONFUNDIR ROBUSTEZ COM PROBABILIDADE
% ============================================================

\begin{frame}{Interpretação da robustez}

    Se uma alternativa ocupar a primeira posição em determinada
    proporção das configurações analisadas, essa medida representa:

    \begin{center}
        \textbf{frequência de primeira posição}
    \end{center}

    e não:

    \begin{center}
        \textbf{probabilidade real de ser a melhor}.
    \end{center}

    \vspace{0.4cm}

    Da mesma forma:

    \[
    \text{robustez}
    \neq
    \text{significância estatística}.
    \]

\end{frame}


% ============================================================
% FRAME 27 — SISTEMA ADAPTATIVO
% ============================================================

\begin{frame}{Por que o sistema é adaptativo?}

    O sistema não será uma tabela fixa de pontuação.

    \vspace{0.3cm}

    Ele deverá responder à estrutura efetivamente observada nos dados:

    \begin{itemize}
        \item critérios disponíveis;
        \item variabilidade;
        \item correlações;
        \item restrições;
        \item alternativas presentes;
        \item condições experimentais.
    \end{itemize}

    \vspace{0.3cm}

    \begin{block}{Adaptação}
        Alterações justificadas nos dados ou critérios provocam a
        reexecução das etapas dependentes, mantendo a análise
        reproduzível.
    \end{block}

\end{frame}


% ============================================================
% FRAME 28 — METODOLOGIA COMPLETA
% ============================================================

\begin{frame}{Síntese metodológica}

    \begin{center}

        Dados

        \[
        \downarrow
        \]

        Auditoria e base canônica

        \[
        \downarrow
        \]

        ADA e seleção dos critérios

        \[
        \downarrow
        \]

        Normalização

        \[
        \downarrow
        \]

        CRITIC

        \[
        \downarrow
        \]

        TOPSIS

        \[
        \downarrow
        \]

        Cenários + Pareto + Sensibilidade

        \[
        \downarrow
        \]

        Robustez e interpretação

    \end{center}

\end{frame}


% ============================================================
% INFORMAÇÕES A SEREM PREENCHIDAS PELOS INTEGRANTES
% ============================================================
%
% PRIMARY_NORMALIZATION = TO_BE_FILLED
% N_CRITERIA = TO_BE_FILLED
% CRITERIA = TO_BE_FILLED
% CRITIC_CORRELATION_METHOD = TO_BE_FILLED
% OFFICIAL_SCENARIOS = TO_BE_FILLED
% WEIGHT_PERTURBATION_RULE = TO_BE_FILLED
% PARETO_TOLERANCE = TO_BE_FILLED
% ROBUSTNESS_METRICS = TO_BE_FILLED
%
% ============================================================


% ============================================================
% MATERIAL EXTERNO / REFERÊNCIAS
% ============================================================
%
% Durante a execução, os integrantes deverão pesquisar e
% confirmar referências metodológicas reais para:
%
% - CRITIC;
% - TOPSIS;
% - demais métodos que necessitem fundamentação.
%
% Não inserir referências inventadas.
%
% Caso a professora forneça metodologia, requisito ou orientação
% adicional, esse material deverá ser incorporado e identificado.
%
% ============================================================


% ============================================================
% ESTADO ATUAL
% ============================================================
%
% METHODOLOGY_SECTION_STATUS = BASE_PREPARED
%
% CRITERIA_STATUS = NOT_DEFINED
% NORMALIZATION_STATUS = NOT_DEFINED
% CRITIC_STATUS = NOT_EXECUTED
% TOPSIS_STATUS = NOT_EXECUTED
% PARETO_STATUS = NOT_EXECUTED
% SENSITIVITY_STATUS = NOT_EXECUTED
% ROBUSTNESS_STATUS = NOT_EXECUTED
%
% ============================================================